\begin{theorem}[Theorem 3]
Let $n\ge2$, $p>1$, $0<\tau<n$, $0<\lambda\le\Lambda$, $a>0$,
$0<\nu_0<\nu_1$, $d>0$, $M\ge0$, $W\ge0$ and let $\sigma\colon\R\to\R$, and suppose the Sobolev--Poincar\'e inequality on balls holds in dimension $n$
(\noderef{CitedSobolevPoincareBall}), and assume the cited regularity result
(C1) (\noderef{CitedRegularityEpsPos}), the cited solvability result (C2)
(\noderef{CitedDirichletW1p}) and the cited Dirichlet result (C3) with H\"older
boundary data (\noderef{CitedDirichletHolder}) in the universal forms described
below.  Then
there are $N_1,N_2\ge0$, depending only on these quantities, with the following
property.

Let $0\le\varepsilon\le1$, let $\Omega\subseteq\R^n$ be open and let
$\Omega'\subseteq\R^n$ be open, nonempty and bounded and a bounded Lipschitz
domain (\noderef{IsBoundedLipschitzDomain}), with
$\overline{\Omega'}\subseteq\Omega$ and $d\le d(\Omega',\partial\Omega)$ --- ``$\Omega'$ is a bounded Lipschitz domain'' being the hypothesis the paper states at l.~1274, now carried as the predicate \noderef{IsBoundedLipschitzDomain} and no longer only as the three clauses open, nonempty and bounded; those three are kept beside it, being implied by it, so that the proof sites discharging them by name read unchanged ---
and assume the cited Campanato embedding (C5) on
$\Omega'$ (\noderef{CitedCampanatoEmbedding}), and also the cited
H\"older estimates (C4) (\noderef{CitedHolderOscillationMinimum}).  Let $f$ satisfy
(F1) of the thesis, (F2) of the thesis and (F3) of the thesis on $\Omega$ with parameters
$p,\varepsilon,\lambda,\Lambda,a,\nu_0,\nu_1$ and modulus $\sigma$
(\noderef{VariationalStructureF123}), and let
$u\in W^{1,p}(\Omega)\cap L^\infty(\Omega)$ with weak gradient $g=\nabla u$
(\noderef{MemW1pOn}), $|u|\le M$ a.e.\ and
$\int_\Omega|u|^p\,dx+\int_\Omega|\nabla u|^p\,dx\le W^p$ (that is,
$\|u\|_{W^{1,p}(\Omega)}\le W$), be a local minimum for $F$ in
$\Omega$ (\noderef{IsLocalMinimumOn}).  Then
\begin{itemize}
\item[(i)] $\|\,|\nabla u|\,\|_{L^{p,n-\tau}(\Omega')}\le N_1$, i.e.\
  $|\nabla u|\in L^{p,n-\tau}(\Omega')$ (\noderef{MorreyBoundOn});
\item[(ii)] $\|u\|_{\mathcal L^{\frac n{n-1},\,n+(1-\frac\tau p)(\frac n{n-1})}(\Omega')}\le N_2$,
  i.e.\ $u\in\mathcal L^{\frac n{n-1},\,n+(1-\frac\tau p)(\frac n{n-1})}(\Omega')$
  (\noderef{CampanatoBoundOn});
\item[(iii)] for each $0\le\eta<1$, $u$ agrees almost everywhere on $\Omega'$
  with a function in $C^\eta(\overline{\Omega'})$ (\noderef{MemHolderOn}), i.e.\
  $u\in C^\eta(\Omega')$.
\end{itemize}
The norms $N_1$, $N_2$ are bound before $\varepsilon$, $\Omega$, $\Omega'$, $f$
and $u$, and after $d(\Omega',\partial\Omega)$, $\sigma$,
$\tau,\lambda,\Lambda,n,p,a,\nu_0,\nu_1$, $\|u\|_{L^\infty(\Omega)}$ and
$\|u\|_{W^{1,p}(\Omega)}$ (here represented by $M$ and $W$): this is the formal
content of the paper's remark that these norms depend exclusively on those
quantities.  No quantity depending on $|\Omega|$ appears.

\noindent\textbf{On the $W^{1,p}$ datum.}  The hypothesis
$\int_\Omega|u|^p\,dx+\int_\Omega|\nabla u|^p\,dx\le W^p$ is a bound on the
genuine Sobolev norm $\|u\|_{W^{1,p}(\Omega)}$, which is the quantity the
paper's dependency list names.  It must not be replaced by a bound on
$\int_\Omega(1+|\nabla u|^2)^{p/2}dx$: that integrand is pointwise at least $1$,
so such a bound would force $W\ge|\Omega|$, would be unsatisfiable whenever
$|\Omega|=\infty$ --- and the paper states this result for an arbitrary open
$\Omega\subseteq\R^n$, so $\Omega=\R^n$ is an instance --- and would let the
conclusion's norm degrade with $|\Omega|$, which is not on the paper's list.
The weighted integral is still available locally where the proof needs it: since
$(1+t^2)^{p/2}\le2^{p/2}(1+t^p)$ for $t\ge0$,
\[
\int_{B_R}(1+|\nabla u|^2)^{\frac p2}dx
\;\le\;2^{\frac p2}\Bigl(|B_R|+\int_{B_R}|\nabla u|^p\,dx\Bigr)
\;\le\;2^{\frac p2}\bigl(c_nR^{\,n}+W^p\bigr),
\]
in which every term is controlled by $n$, $p$, $R$ and $W$ alone.

\medskip
\noindent\textbf{Correction reported.}  Following the paper's Note for the
formalization, $\tau$ ranges over $(0,n)$ and not $[0,n)$.  For $\tau=0$ part
(i) would assert $|\nabla u|\in L^{p,n}=L^\infty$, and the proof needs $\tau>0$
in order to apply \noderef{CampanatoIterationLemma} with
$\beta=n-\tau<\alpha=n$.

The hypothesis \noderef{CitedSobolevPoincareBall} is the cited result (C6) of the
paper's Note for the formalization, carried here as an explicit argument because
the proof uses it --- it enters through the constant $c_0$ of
\noderef{FrozenComparisonDecay} and through the constants $N_1,N_2$ of
\noderef{ComparisonSchemeMorreyCore}, which those two nodes build with it ---
which is what the Note permits for the results the thesis quotes without proof.  Although it is
bound before the existential for $N_1$ and $N_2$, it introduces no dependence of
$N_1$ and $N_2$ on any quantity outside the list the paper names: (C6) is a statement
about $n$ alone, and $n$ is already bound first, so the hypothesis carries no
numerical content of its own.

\medskip
\noindent\textbf{The three cited results for the frozen operator.}  The cited
regularity result (C1) (\noderef{CitedRegularityEpsPos}), the cited solvability
result (C2) (\noderef{CitedDirichletW1p}) and the cited Dirichlet result (C3)
with H\"older boundary data (\noderef{CitedDirichletHolder}) are assumed here
for \emph{every} field $A'\colon\R^n\to\R^n$ and all parameters
$\varepsilon',\gamma_0',\gamma_1'$ with $\gamma_0'>0$ and $\gamma_1'>0$ --- and,
for (C1), for every open $\Omega_0\subseteq\R^n$ as well --- which are the forms
these three results have in the paper's Note for the formalization, where they
``may be taken as explicit hypotheses \ldots\ where they are used''.

The universal shape is forced rather than convenient.  The operator to which all
three are applied is the frozen operator $A=\nabla_hf^0$ with
$f^0(h)=f(x_0,u(x_0),h)$ (\noderef{FrozenGradientOperator}), which varies with
the centre $x_0$ and with the value $u(x_0)$; both are chosen inside the proof,
after the existential for $N_1$ and $N_2$, so no single instance of (C1)--(C3)
can be named in the statement.  All three are, in this same universal form, hypotheses of
\noderef{FrozenComparisonDecay}, and the proof reaches all three through that
node: (C1) and (C2) because that node obtains the decay (4.11) of the thesis from
Theorem~1, and (C3) because that node produces the comparison map.  Like (C6), none of the three introduces a
dependence of $N_1,N_2$ on any quantity outside the list the paper names: each
mentions only $n$ and $p$, both already bound, and each is universally
quantified over its remaining parameters, so none carries numerical content of
its own.  The guards $\gamma_0'>0$, $\gamma_1'>0$ are the thesis' own standing
positivity conditions on the structure constants, and the openness guard on (C1)
is the domain condition \noderef{CitedRegularityEpsPos} is stated for; (C2) and
(C3) need no such guard because \noderef{CitedDirichletW1p} and
\noderef{CitedDirichletHolder} are statements about balls and fix their own
domains.
\textbf{On non-vacuity of the structure hypothesis.}  The statement is
conditioned on \noderef{VariationalStructureF123}, so it would be vacuously true
if no density satisfied that bundle, which the Note for the formalization forbids
at l.~258.  It does not: \noderef{ModelDensityStructureF123} exhibits the density
$f(h)=(\varepsilon+|h|^2)^{p/2}$ of \S4, with $\sigma\equiv0$, as a member of the
bundle for every $\varepsilon\ge0$, every $a\ge0$ and every $p$ in the singular
range $1<p<2$, at $\lambda=\Lambda=1$, $\nu_0=p(p-1)$ and $\nu_1=p$ --- constants
admissible against the hypotheses $0<\lambda\le\Lambda$ and $0<\nu_0<\nu_1$
imposed here.  It also satisfies the extra condition (F4) of the thesis that
\noderef{MinimumGradientHolder} adds, at every $c_\sigma>0$ and $\delta>0$.  The
citation is for non-vacuity only; no step of the proof below uses the model
density.
\end{theorem}
\begin{proof}
Throughout write
\[
q=\frac n{n-1},\qquad
\theta=n+q\Bigl(1-\frac\tau p\Bigr)
=n+\Bigl(1-\frac\tau p\Bigr)\frac n{n-1},
\qquad
\omega_n=|B_1(0)|\in(0,\infty)
\]
(the Preamble item on the volume of a ball), so that $\theta$ is the Campanato
exponent appearing in (ii) and $q>1$.

The proof follows the thesis' own division of labour at p.~50.  What is
specific to the variational setting is the verification that the frozen
operator satisfies (Q1) of the thesis--(Q2) of the thesis, which is read off from
(F2) of the thesis through \noderef{FrozenOperatorStructureQ12} in Step 3, and the
smallness of the remainder, which Lemma~4.1 supplies in Step 4.  The first of
those two is itself a node, \noderef{FrozenOperatorStructureQ12}, because \S3's
statements are all posed for a field satisfying
(Q1) of the thesis--(Q2) of the thesis and three of them consume it.  Everything else is
carried by three nodes that serve this theorem and Theorem~5 alike: \noderef{InnerDomainRepresentative} produces the
inner domain and the continuous representative (Step 2),
\noderef{FrozenComparisonDecay} produces the comparison map of (4.1) of the thesis
and the decay (4.11) of the thesis of its gradient (Step 3), and
\noderef{ComparisonSchemeMorreyCore} carries the whole passage from
(4.11) of the thesis and $(\star)$ to the three conclusions --- the Morrey decay
(4.14) of the thesis, the Morrey bound (i), the Campanato bound (ii) and the
H\"older representative (iii) --- in Step 5.

\emph{Step 1: the constants.}  They are all chosen now, before
$\varepsilon,\Omega,\Omega',f,u,g$, which is what the quantifier order of the
statement demands.

\begin{enumerate}
\item[(1a)] Put $r=\nu_1/\nu_0>0$ and apply \noderef{FrozenComparisonDecay}
  with the parameters $n,p$ and this $r$: its hypotheses $2\le n$, $1<p$,
  $0<r$ hold, its hypothesis \noderef{CitedSobolevPoincareBall} is the present
  theorem's, and its three universal hypotheses --- the cited solvability
  result (C2) (\noderef{CitedDirichletW1p}), the cited regularity result (C1)
  (\noderef{CitedRegularityEpsPos}) and the cited Dirichlet result (C3) with
  H\"older boundary data (\noderef{CitedDirichletHolder}) --- are the present
  theorem's (C2), (C1) and (C3) clauses verbatim, in the same universal forms:
  (C2) and (C3) for every field $A'$ and all $\varepsilon',\gamma_0',\gamma_1'$
  with $\gamma_0'>0$ and $\gamma_1'>0$, and (C1) for all of those \emph{and} for
  every open set $\Omega_0$, that clause alone carrying a set in its binder list
  because (C1) is a statement about weak solutions on an open set.  Let $c_0>0$ be the constant
  it produces; $c_0$ depends only on $n,p$ and $\nu_1/\nu_0$.
\item[(1b)] Apply \noderef{ComparisonGradientEstimateMin} (Lemma~4.1) with the
  parameters $n,p,\lambda,\Lambda,a,\nu_0,\nu_1$.  That node binds five
  hypotheses before the two clauses of its conclusion, and all five are
  hypotheses of the present theorem: $2\le n$, $1<p$,
  $0<\lambda\le\Lambda$, $0<a$ and $0<\nu_0<\nu_1$; the Sobolev--Poincar\'e
  inequality on balls in dimension $n$ (\noderef{CitedSobolevPoincareBall});
  and the cited Dirichlet result (C3) with H\"older boundary data
  (\noderef{CitedDirichletHolder}) for every field $A'$ and all
  $\varepsilon',\gamma_0',\gamma_1'$ with $\gamma_0'>0$ and $\gamma_1'>0$,
  which is the present theorem's (C3) clause verbatim.  Of the two
  clauses of its conclusion take clause (A), which is Lemma 4.1 itself; the
  present theorem assumes no rate for $\sigma$, so clause (B) is not invoked.
  Read clause (A) at $M$, at $d/2$ in the role of its $d$, and at the present
  $\sigma$, whose guards $0\le M$ and $0<d/2$ hold.  Let $\sigma_2$ and
  $R_1>0$ be the function and the radius it produces, so $\sigma_2$ is
  nonnegative, bounded, continuous and increasing with $\sigma_2(0)=0$, and
  both depend only on $n,p,\lambda,\Lambda,a,\nu_0,\nu_1,M,d$ and $\sigma$.
\item[(1c)] Put
  \[
  R_\ast=\min\Bigl\{R_1,\ \frac d{16}\Bigr\}>0
  \]
  and apply \noderef{ComparisonSchemeMorreyCore} with the parameters $n,p,\tau$,
  with $c_0$ of (1a), with $W$ and $M$, with $d/2$ in the role of its $d$, with
  $R_\ast$ in the role of its $R_1$, and with $\sigma_2$ of (1b): its hypotheses
  $2\le n$, $1<p$, $0<\tau<n$, $0<c_0$, $0\le W$, $0\le M$, $0<d/2$,
  $0<R_\ast$, the four properties of $\sigma_2$ (nonnegativity, continuity,
  monotonicity and $\sigma_2(0)=0$) and the Sobolev--Poincar\'e inequality
  (\noderef{CitedSobolevPoincareBall}) hold by (1a), (1b) and the hypotheses of
  the present theorem.  Let $N_1,N_2\ge0$ be the two constants it produces.
  They depend only on $n,p,\tau,c_0,W,M,d,R_\ast$ and $\sigma_2$, hence --- by
  (1a) and (1b) --- only on
  $n,p,\tau,\lambda,\Lambda,a,\nu_0,\nu_1,d,M,W$ and $\sigma$, which is the
  list the statement names, and in particular not on
  $\varepsilon,\Omega,\Omega'$, the data of the problem or $u$ and $g$.
\end{enumerate}

Now fix $\varepsilon\in[0,1]$, the open sets $\Omega,\Omega'$, the density $f$,
the function $u$ and the field $g$ satisfying all the hypotheses.

\emph{Step 2: the inner domain and the continuous representative.}  Apply the
cited H\"older estimates (C4)(a)
(\noderef{CitedHolderOscillationMinimum}) with the parameters
$\lambda,\Lambda,a,M$: its hypotheses $0<\lambda$, $\lambda\le\Lambda$,
$0<a$, $0\le M$ are hypotheses of the present theorem.  Let $\alpha\in(0,1]$
and $d_0>0$ be the constants it produces; they depend only on
$n,p,\lambda,\Lambda,a$ and $M$, and in particular not on $\varepsilon$, which
that clause binds after them.  Read its conclusion at the present
$\varepsilon$, whose guards $0\le\varepsilon$ and $\varepsilon\le1$ are
hypotheses of the present theorem.  It applies to the present data: $\Omega$ is
open; the growth condition (F1) of the thesis with parameters
$p,\varepsilon,\lambda,\Lambda,a$ on $\Omega$ (\noderef{VariationalGrowthF1})
consists of exactly the three clauses (measurability in $x$, continuity in
$(z,h)$ and the two-sided growth bound) that are the first three clauses of
(F1) of the thesis--(F3) of the thesis (\noderef{VariationalStructureF123}), which the
present theorem assumes; $u$ is a local minimum with weak gradient $g$
(\noderef{IsLocalMinimumOn}); and $|u|\le M$ almost everywhere on $\Omega$.  So
for every $y$ and $R_3$ with $\overline B_{R_3}(y)\subseteq\Omega$ and every
$0<\rho\le R_3$,
\[
\operatorname*{osc}\bigl(u,B_\rho(y)\bigr)\le4^\alpha
\max\bigl\{\operatorname*{osc}\bigl(u,B_{R_3}(y)\bigr),\ d_0R_3\bigr\}
\Bigl(\frac\rho{R_3}\Bigr)^\alpha
\tag{1}
\]
(\noderef{EssOscOn}).

Apply \noderef{InnerDomainRepresentative} with the parameters
$n,p,\alpha,d_0,M,d$ --- its hypotheses $1\le n$, $1\le p$ (because $p>1$),
$0<\alpha$, $\alpha\le1$, $0<d_0$, $0\le M$ and $0<d$ hold --- and with the
data $\Omega,\Omega',u,g$: the set $\Omega$ is open, $\Omega'$ is nonempty and
bounded, $d\le|x-y|$ for every $x\in\Omega'$ and every $y\notin\Omega$,
$u\in W^{1,p}(\Omega)$ with weak gradient $g$ (\noderef{MemW1pOn}),
$|u|\le M$ almost everywhere on $\Omega$, and the oscillation decay is $(1)$.
Let $\Omega_1$ and $U$ be the open set and the function it produces, so that
\begin{itemize}
\item[(2a)] $\Omega_1$ is open and bounded,
  $\overline{\Omega'}\subseteq\Omega_1$ and
  $\overline{\Omega_1}\subseteq\Omega$; in particular
  $\Omega_1\subseteq\Omega$;
\item[(2b)] $\frac d2\le|x-y|$ for every $x\in\Omega'$ and every
  $y\notin\Omega_1$, and $\frac d2\le|x-y|$ for every $x\in\Omega_1$ and every
  $y\notin\Omega$;
\item[(2c)] $\overline B_{d/2}(x)\subseteq\Omega$ for every $x\in\Omega_1$;
\item[(2d)] $U=u$ almost everywhere on $\Omega$, and $U\in W^{1,p}(\Omega)$
  with weak gradient $g$ (\noderef{MemW1pOn});
\item[(2e)] $U$ is continuous on $\Omega_1$, and
  $U\in C^\alpha\bigl(\overline B_{d/16}(x)\bigr)$ (\noderef{MemHolderOn}) for
  every $x\in\Omega_1$.
\end{itemize}

Two consequences of $(2d)$.  First, $|U|\le M$ almost everywhere on $\Omega$,
because $U=u$ almost everywhere there.  Second, $U$ is a local minimum for $F$
in $\Omega$ with weak gradient $g$ (\noderef{IsLocalMinimumOn}): apply
\noderef{AEModificationInvariance} with the parameters $n,p$ --- its hypotheses
$1\le n$, $1\le p$ hold --- on the open set $\Omega$, with $u$ in the role of
its $f$, $U$ in the role of its $f'$ and the present $g$, the
almost-everywhere equality being that of $(2d)$; its clause (5), taken at the
present density $f$, transports the local-minimum property of $u$ to $U$.

\emph{Step 3: the comparison map and the decay (4.11) of the thesis.}  Fix
$x_1\in\Omega_1$ and $0<R\le R_\ast$; note $R\le R_\ast\le d/16$, so
$\overline B_R(x_1)\subseteq\overline B_{d/2}(x_1)\subseteq\Omega$ by $(2c)$.
Put
\[
z_1=U(x_1),\qquad f^0(h)=f(x_1,z_1,h),\qquad
A^0=\nabla_hf^0\quad(\noderef{FrozenGradientOperator}),\qquad
\widetilde A=A^0-A^0(0).
\]
The normalised frozen operator $\widetilde A$ satisfies
(Q1) of the thesis--(Q2) of the thesis on all of $\R^n$ with the parameters
$p,\varepsilon$ and the constants $\gamma_0=\nu_0$, $\gamma_1=\nu_1$
(\noderef{EllipticStructureQ12}): this is \noderef{FrozenOperatorStructureQ12},
applied with the parameters $n,p,\varepsilon,\lambda,\Lambda,a,\nu_0,\nu_1$ ---
whose guards $2\le n$, $1<p$, $0\le\varepsilon\le1$, $0<\lambda\le\Lambda$,
$0<a$, $0<\nu_0<\nu_1$ are hypotheses of the present theorem --- and with the
data $\Omega$ (open), the density $f$, the modulus $\sigma$, the point $x_1$ and
the value $z_1$: its remaining hypotheses are $x_1\in\Omega$, which holds
because $x_1\in\Omega_1\subseteq\Omega$ by $(2a)$, and
(F1) of the thesis--(F3) of the thesis on $\Omega$ with those parameters and that modulus
(\noderef{VariationalStructureF123}), which is a hypothesis of the present
theorem.

Now apply the second part of \noderef{FrozenComparisonDecay} --- the same
instance as in (1a), so the constant it produces is the $c_0$ fixed there ---
with the data $\varepsilon$, $\alpha$, $\gamma_0=\nu_0$, $\gamma_1=\nu_1$, the
field $A^0$, the centre $x_1$, the radius $R$, the function $U$ and the field
$g$.  Its hypotheses hold:
\begin{itemize}
\item $0\le\varepsilon$ and $\varepsilon\le1$; $0<\alpha$ and $\alpha\le1$
  (Step 2); $0<\nu_0$ and $0<\nu_1$; $\nu_1\le r\nu_0$, with equality by the
  definition of $r$ in (1a); and $0<R$;
\item $\widetilde A=A^0-A^0(0)$ satisfies (Q1) of the thesis--(Q2) of the thesis on
  $\R^n$ with the parameters $p,\varepsilon,\nu_0,\nu_1$
  (\noderef{EllipticStructureQ12}), which is the previous paragraph;
\item $U\in W^{1,p}(B_R(x_1))$ with weak gradient $g$ (\noderef{MemW1pOn}), by
  $(2d)$ and clause (1) of \noderef{WeakGradientRestrict} applied with $q=p$ to
  the open subset $B_R(x_1)\subseteq\Omega$: its outer hypotheses are $1\le n$
  (from $2\le n$), $1\le p$ (from $1<p$), the openness of $\Omega$ and of
  $B_R(x_1)$, and the inclusion $B_R(x_1)\subseteq\Omega$, which holds because
  $B_R(x_1)\subseteq\overline B_R(x_1)\subseteq\Omega$ as just noted, and its
  clause-(1) hypothesis is the $U\in W^{1,p}(\Omega)$ of $(2d)$;
\item $U\in C^\alpha\bigl(\overline B_R(x_1)\bigr)$ (\noderef{MemHolderOn}):
  by $(2e)$ at the centre $x_1\in\Omega_1$ we have
  $U\in C^\alpha(\overline B_{d/16}(x_1))$, and
  $\overline B_R(x_1)\subseteq\overline B_{d/16}(x_1)$ because
  $R\le R_\ast\le d/16$; a function of class $C^\alpha$ on a set is of class
  $C^\alpha$ on every subset of it, the continuity restricting and the same
  nonnegative constant serving.
\end{itemize}
It produces functions $v,V$ and a vector field $g_v$ such that the bundle
$\mathrm{FrozenComparisonOn}(n,p,x_1,R,A^0,U,U,g,v,V,g_v)$ holds
(\noderef{FrozenComparisonOn}) --- so in particular $v$ is a weak solution of
$\operatorname{div}A^0(\nabla v)=0$ in $B_R(x_1)$
(\noderef{IsWeakSolutionOn}), $v\in W^{1,p}(B_R(x_1))$ with weak gradient $g_v$
(\noderef{MemW1pOn}), and $v-U\in W^{1,p}_0(B_R(x_1))$ with weak gradient
$g_v-g$ (\noderef{MemW1pZeroOn}) --- and such that
\[
\int_{B_\rho(x_1)}\bigl(\varepsilon+|g_v|^2\bigr)^{\frac p2}dx
\le c_0\Bigl(\frac\rho R\Bigr)^{n}
\int_{B_R(x_1)}\bigl(\varepsilon+|g_v|^2\bigr)^{\frac p2}dx
\qquad(0<\rho\le R),
\tag{4.11}
\]
which is the first of the two estimates the comparison scheme asks for.

\emph{Step 4: the smallness $(\star)$.}  Apply
\noderef{ComparisonGradientEstimateMin} (Lemma~4.1) with the parameters of (1b)
--- so clause (A) of its conclusion, with $d/2$ in the role of its $d$,
and with the $\sigma_2$ and $R_1$ it produced --- and with the data
$\varepsilon$, $\Omega$, $\Omega_1$ in the role of its $\Omega'$,
$f$, $U$ in the role of its $u$ \emph{and} of its $U$, the field $g$, the
centre $x_1$, the radius $R$, and the $v,V,g_v$ of Step 3.  Its hypotheses
hold: $0\le\varepsilon\le1$; $\Omega$ is open; $\Omega_1$ is open and
bounded with $\overline{\Omega_1}\subseteq\Omega$ and, by $(2b)$,
$d/2\le|x-y|$ for $x\in\Omega_1$, $y\notin\Omega$;
(F1) of the thesis--(F3) of the thesis hold on $\Omega$ with the stated parameters
(\noderef{VariationalStructureF123}); the cited H\"older estimates (C4) are
assumed (\noderef{CitedHolderOscillationMinimum}); $U\in W^{1,p}(\Omega)$ with
weak gradient $g$ and $|U|\le M$ almost everywhere on $\Omega$ and
$U$ is a local minimum for $F$ in $\Omega$, all three by Step 2;
$x_1\in\Omega_1$; $0<R\le R_\ast\le R_1$;
$\overline B_R(x_1)\subseteq\Omega$ (Step 3); and the frozen comparison bundle
of Step 3, whose frozen operator is $\nabla_hf^0$ with
$f^0(h)=f(x_1,U(x_1),h)$, as that node requires.  Its conclusion is
\[
\int_{B_R(x_1)}|g-g_v|^p\,dx
\le\sigma_2(R)\int_{B_R(x_1)}(1+|g|^2)^{\frac p2}dx .
\tag{$\star$}
\]

\emph{Step 5: the comparison scheme, and the three conclusions.}  Apply
\noderef{ComparisonSchemeMorreyCore} with the parameters of (1c) --- so with
$n,p,\tau$, with $c_0$ of (1a), with $W$ and $M$, with $d/2$ in the role of its
$d$, with $R_\ast$ in the role of its $R_1$ and with $\sigma_2$ of (1b) --- and
with the data $\varepsilon$, $\Omega$, $\Omega_1$ in the role of
its $\Omega''$, $\Omega'$, the function $u$ and the field $g$.  Its hypotheses
hold:
\begin{itemize}
\item $0\le\varepsilon\le1$; $\Omega$, $\Omega_1$ and $\Omega'$ are
  open, $\Omega'$ is nonempty by hypothesis and $\Omega_1$ is bounded, and
  $\overline{\Omega'}\subseteq\Omega_1$ and
  $\overline{\Omega_1}\subseteq\Omega$, all by $(2a)$;
\item the two separations, both with the constant $d/2$, which are $(2b)$;
\item $\Omega'$ a bounded Lipschitz domain
  (\noderef{IsBoundedLipschitzDomain}), a hypothesis of the present theorem,
  which that node now carries because the (C5) instance it assumes requires it;
\item the cited Campanato embedding (C5) on $\Omega'$ at the exponent
  $n/(n-1)$ (\noderef{CitedCampanatoEmbedding}) and the Sobolev--Poincar\'e
  inequality on balls (\noderef{CitedSobolevPoincareBall}) are hypotheses of the
  present theorem;
\item $u\in W^{1,p}(\Omega)$ with weak gradient $g$ (\noderef{MemW1pOn}),
  $|u|\le M$ almost everywhere on $\Omega$, and
  $\int_\Omega|u|^p\,dx+\int_\Omega|\nabla u|^p\,dx\le W^p$: these are
  hypotheses of the present theorem;
\item the comparison bundle, which is where Steps 2--4 are consumed: the
  centre $x_1\in\Omega_1$ and the radius $0<R\le R_\ast$ were arbitrary in
  Steps 3--4, so for every such pair the field $g_v$ produced in
  Step 3 satisfies (4.11) of the thesis by Step 3 and $(\star)$ by Step 4 --- and
  $\overline B_R(x_1)\subseteq\Omega$ holds automatically for such a pair, by
  $(2c)$ and $R_\ast\le d/16$, so the node's guard
  $\overline B_R(x_1)\subseteq\Omega$ is met.  Note that $R_\ast\le R_1$,
  which is what Step 4 requires, and $R_\ast\le d/16$, which is what Step 3
  requires.
\end{itemize}
Its three conclusions, with the $N_1$ and $N_2$ of (1c), are verbatim the three
assertions of the present theorem: the Morrey bound
$\|\,|\nabla u|\,\|_{L^{p,n-\tau}(\Omega')}\le N_1$
(\noderef{MorreyBoundOn}), which is (i); the Campanato bound
$\|u\|_{\mathcal L^{q,\theta}(\Omega')}\le N_2$ with
$\theta=n+(1-\frac\tau p)\frac n{n-1}$ (\noderef{CampanatoBoundOn}), which is
(ii); and, for each $\eta\in[0,1)$, a function $w$ equal to $u$ almost
everywhere on $\Omega'$ and of class $C^\eta(\overline{\Omega'})$
(\noderef{MemHolderOn}), which is (iii).  This completes the proof.
\end{proof}
